% Continuous Gradient-Energy Importance Accumulation Reduces Catastrophic Forgetting
% ContinualAI Workshop @ NeurIPS 2026 submission
%
% TRACEABILITY NOTE: All numerical results traced to
%   tar_state/honest_evidence_inventory.json
%   tar_state/comparisons/phase10_controlled_rerun_20260509T132155Z.json
%   tar_state/comparisons/phase11_ablation__20260511T113318Z.json
%
% For final submission: replace \documentclass[...]{article} with
%   \usepackage{neurips_2024} per workshop instructions.
%
% Compile: pdflatex main.tex && bibtex main && pdflatex main.tex && pdflatex main.tex

\documentclass[10pt,letterpaper]{article}

% ── packages ──────────────────────────────────────────────────────────────────
\usepackage[margin=1in]{geometry}
\usepackage{amsmath,amssymb,amsthm}
\usepackage{booktabs}
\usepackage{multirow}
\usepackage{xcolor}
% \usepackage{microtype}  % disabled: requires scalable font setup
\usepackage[numbers,sort&compress]{natbib}
\usepackage[colorlinks=true,linkcolor=blue,citecolor=blue,urlcolor=blue]{hyperref}

% ── theorem environments ──────────────────────────────────────────────────────
\newtheorem{theorem}{Theorem}
\newtheorem{corollary}[theorem]{Corollary}
\newtheorem{proposition}[theorem]{Proposition}
\newtheorem{remark}{Remark}

% ── custom macros ─────────────────────────────────────────────────────────────
\newcommand{\TCL}{\textsc{TCL}}
\newcommand{\EWC}{\textsc{EWC}}
\newcommand{\SI}{\textsc{SI}}
\newcommand{\SGD}{\textsc{SGD}}
\newcommand{\forgetmean}[1]{\ensuremath{#1_{\mathrm{fgt}}}}
\newcommand{\accmean}[1]{\ensuremath{#1_{\mathrm{acc}}}}
\newcommand{\CI}[2]{[\,{#1},\,{#2}\,]}

% ── title / authors ───────────────────────────────────────────────────────────
\title{Continuous Gradient-Energy Importance Accumulation\\
Reduces Catastrophic Forgetting:\\
An Empirical Audit of \EWC{}'s Fisher Snapshot Assumption}

\author{
  Anonymous Author(s)\\
  \texttt{anonymous@institution.edu}
}

\date{ContinualAI Workshop @ NeurIPS 2026}

% ══════════════════════════════════════════════════════════════════════════════
\begin{document}
\maketitle

% ── Abstract ──────────────────────────────────────────────────────────────────
\begin{abstract}
Elastic Weight Consolidation (\EWC) estimates per-parameter importance as a
uniform average of squared gradients from the final training batches of each
task---a Fisher snapshot that conflates high-noise early-training gradients
with the convergence-phase curvature that Fisher information actually captures.
We propose \emph{continuous gradient-energy importance accumulation}: an
exponential moving average (EMA) of squared gradients maintained throughout
training, which we show converges to the diagonal Fisher information with
exponentially decaying bias while down-weighting early-training noise.
The resulting regularizer, \TCL{}, applies an elastic penalty using EMA-derived
importance weights in place of \EWC{}'s terminal snapshot.

On Split-CIFAR-10 (ResNet-18, 5 tasks, $n{=}5$ seeds, 40 epochs/task), \TCL{}
achieves mean catastrophic forgetting $0.116$
($95\%$~CI $\CI{0.080}{0.153}$) versus $\SGD$'s $0.233$
($\CI{0.229}{0.237}$), a reduction of $0.117$
($p{=}0.0012$, Cohen's $d{=}3.68$, $5/5$ seeds, Bonferroni-corrected $k{=}4$).
Against \EWC{} ($\lambda{=}100$), \TCL{} shows directional improvement
$\Delta{=}{-}0.075$ ($p{=}0.019$, $d{=}1.70$, $5/5$ seeds),
which does not survive Bonferroni correction
($\alpha/k{=}0.0125$, $p{=}0.019 > 0.0125$).

A pre-registered mechanism ablation definitively establishes that the elastic
penalty is the sole productive component: a governor-alone condition (regime-detection
without penalty) performs \emph{worse} than plain \SGD{} ($0.250$ vs.\ $0.219$),
and the thermodynamic learning-rate adjustment never activates in any production trace
($\rho{=}0$, $\sigma{=}0$).
We conclude that continuous importance accumulation, not thermodynamic regime detection,
explains the empirical benefit.
\end{abstract}

% ── 1. Introduction ───────────────────────────────────────────────────────────
\section{Introduction}

Neural networks trained sequentially on a stream of tasks suffer from
\emph{catastrophic forgetting}: optimizing for a new task overwrites parameters
critical for previously learned tasks~\cite{McCloskey1989}.
Regularization-based continual learning methods address this by penalizing changes
to parameters deemed important for past tasks, with the importance estimated from
the Fisher information matrix~\cite{Kirkpatrick2017} or path integrals of parameter
movements~\cite{Zenke2017}.

The dominant approach, \EWC{}~\cite{Kirkpatrick2017}, approximates the diagonal
Fisher by averaging squared gradients over the final $N$ training batches
immediately after each task completes.
This snapshot has two structural weaknesses.
First, gradients early in training reflect random initialization and high loss
rather than the task's true parameter importance; mixing these with late-training
gradients biases the importance estimate upward for parameters with large early-training
noise.
Second, the snapshot is computed \emph{once} from a terminal window and never
updated as training continues on future tasks, discarding information from the
entire non-terminal training trajectory.

\paragraph{This paper.}
We propose replacing \EWC{}'s terminal snapshot with a continuous EMA of
squared gradients---the \emph{ThermalImportance} accumulator.
We show (Theorem~\ref{thm:ema_convergence}) that this accumulator converges to
the diagonal Fisher information with exponentially decaying bias, and that
\EWC{}'s snapshot estimator is the limiting case $\beta \to 1$.
Under non-stationary gradient processes (as occur in all neural network training),
the EMA's recency weighting provides a strictly smaller non-stationarity bias
than \EWC{}'s uniform average (Remark~\ref{rem:nonstationarity}).

We evaluate the resulting regularizer, \TCL{},\footnote{%
  TCL originally stood for \emph{Thermodynamic Continual Learning}, reflecting a
  physical intuition that high-gradient parameters are ``hot'' and should be
  stabilized.  Mechanism ablation (Section~\ref{sec:experiments}) shows the
  thermodynamic regime-detection component does not contribute to performance.
  We retain the acronym but clarify that the operative mechanism is purely
  gradient-energy elastic regularization.}
against \EWC{}, \SI{}, and \SGD{} on Split-CIFAR-10 using a pre-registered
controlled experiment with $n{=}5$ random seeds and corrected statistical
methodology (t-distribution confidence intervals, Bonferroni family-wise
correction, non-central $t$ power analysis).

Our contributions are:
\begin{enumerate}
  \item \textbf{Formal analysis.}
    We prove that the ThermalImportance EMA is an online, continuously-updated
    approximation to the diagonal Fisher information, with \EWC{} as the $\beta{\to}1$
    limiting case and exponentially smaller non-stationarity bias (Section~\ref{sec:method}).
  \item \textbf{Controlled empirical comparison.}
    On Split-CIFAR-10 with corrected statistical methodology, \TCL{} reduces
    forgetting by $0.117$ relative to \SGD{} ($p{=}0.0012$, $d{=}3.68$,
    Bonferroni-significant at $k{=}4$).
    Relative to \EWC{}, the improvement is directional ($p{=}0.019$, $d{=}1.70$)
    but does not survive family-wise correction (Section~\ref{sec:experiments}).
  \item \textbf{Mechanism ablation.}
    A pre-registered four-condition ablation establishes that the elastic penalty is
    the sole productive mechanism: governor-alone performs worse than \SGD{},
    and the LR adjustment never fires in production (Section~\ref{sec:experiments}).
\end{enumerate}

% ── 2. Related Work ───────────────────────────────────────────────────────────
\section{Related Work}
\label{sec:related}

\paragraph{Catastrophic forgetting and evaluation scenarios.}
Sequential learning in artificial neural networks induces catastrophic
interference~\cite{McCloskey1989}: gradient-based updates that minimize a new
task's loss overwrite weight configurations that supported prior tasks.
Van de Ven and Tolias~\cite{vandeVen2019} identify three evaluation scenarios:
task-incremental (task identity provided at test time),
class-incremental (task identity withheld), and domain-incremental.
This paper addresses the \emph{task-incremental} scenario exclusively.

\paragraph{Regularization-based methods.}
\EWC{}~\cite{Kirkpatrick2017} adds an elastic term to the loss function that
penalizes changes to parameters with high Fisher information:
$\mathcal{L}_{\EWC} = \mathcal{L}_{\mathrm{CE}} + \frac{\lambda}{2}\sum_i \hat{F}_{ii}(\theta_i - \theta_i^*)^2$,
where $\hat{F}_{ii}$ is the diagonal of the empirical Fisher estimated from a
post-task snapshot.
\EWC{}'s importance estimate is a uniform average of squared gradients over the
final $N$ batches, discarding the full training trajectory.

Synaptic Intelligence (\SI{})~\cite{Zenke2017} accumulates importance online via
path integrals of parameter movements during training, providing a stronger prior
over slow-moving parameters.
However, \SI{} at published default hyperparameters collapses to chance-level
accuracy on our benchmark (Section~\ref{sec:experiments}), consistent with known
sensitivity to the regularization strength $c$.

Memory Aware Synapses~\cite{Aljundi2018} estimates importance from the gradient
of the output (rather than the loss), reducing dependency on label information.
All three methods share the structural assumption that importance can be reliably
estimated from a finite sample; our work argues that for EWC specifically,
continuous accumulation is more principled.

\paragraph{Replay-based methods.}
Replay methods maintain a memory buffer of past examples~\cite{Rebuffi2017,LopezPaz2017}
or generate synthetic past data~\cite{Shin2017}.
Dark Experience Replay (DER++)~\cite{Buzzega2020} stores soft logits alongside raw
inputs to better preserve the model's decision boundary.
The current work does not compare against replay methods; DER++ and LwF baselines
are included in the pre-registered Phase 2 replication plan but have not yet been run.

\paragraph{Fisher information approximations.}
The diagonal Fisher information $F_{ii} = \mathbb{E}_{x}[(\partial \mathcal{L}/\partial\theta_i)^2]$
is the cornerstone of \EWC{}'s importance weighting.
Efficient diagonal approximations and Kronecker-factored estimates have been
studied for second-order optimization~\cite{Martens2015}, but continual learning
has relied primarily on diagonal estimates computed from task-end snapshots.
Our contribution is to show that replacing the snapshot with a continuously maintained
EMA yields a formally better approximation while remaining computationally identical
in per-parameter cost.

% ── 3. Method ─────────────────────────────────────────────────────────────────
\section{Method}
\label{sec:method}

\subsection{Background: EWC's Fisher Snapshot}
\label{sec:ewc_background}

Let $\theta \in \mathbb{R}^d$ denote network parameters after training on a task,
and let $\mathcal{D}_{\mathrm{task}}$ be the task data distribution.
The diagonal Fisher information for parameter $i$ is
$F_{ii} = \mathbb{E}_{x \sim \mathcal{D}_{\mathrm{task}}}[g_i^2]$,
where $g_i = \partial\mathcal{L}(\theta; x)/\partial\theta_i$.

\EWC{} approximates $F_{ii}$ by the empirical mean over the final $N$ training batches:
\begin{equation}
  \hat{F}_{ii}^{\EWC} = \frac{1}{N}\sum_{t=T-N+1}^{T} g_i(t)^2.
  \label{eq:ewc_snapshot}
\end{equation}
This estimate is unbiased only if gradients are stationary throughout training.
In practice, gradients decay substantially as training converges: early-training
gradients reflect large loss curvature at random initialization, while late-training
gradients reflect the convergence-phase curvature that Fisher information captures.
A uniform window containing both phases overestimates importance for parameters
with large early-training noise.

\subsection{ThermalImportance EMA Accumulator}
\label{sec:thermal_importance}

\TCL{} maintains a running EMA of squared gradients throughout training.
For parameter $i$ at training step $T$:
\begin{equation}
  v_i^{(T)} = (1-\beta)\sum_{t=1}^{T}\beta^{T-t}\,g_i(t)^2,
  \label{eq:ema_def}
\end{equation}
where $\beta \in (0,1)$ is the EMA decay constant.
This accumulator is updated continuously as training proceeds, requiring no
designated importance-extraction phase after task completion.

\begin{theorem}[EMA convergence to diagonal Fisher]
\label{thm:ema_convergence}
Let $g_i(t) \stackrel{\mathrm{i.i.d.}}{\sim} \mathcal{D}_{g_i}$ with
$\mathbb{E}[g_i(t)^2] = F_{ii} < \infty$.
For the EMA accumulator \eqref{eq:ema_def}:
\begin{equation}
  \mathbb{E}[v_i^{(T)}] = F_{ii}(1 - \beta^T).
  \label{eq:bias}
\end{equation}
The bias $|\mathbb{E}[v_i^{(T)}] - F_{ii}| = F_{ii}\,\beta^T \to 0$ exponentially as $T \to \infty$.
\end{theorem}
\begin{proof}
By linearity of expectation,
$\mathbb{E}[v_i^{(T)}] = (1-\beta)\sum_{t=1}^T \beta^{T-t} F_{ii}
 = F_{ii}(1-\beta)\cdot\frac{1-\beta^T}{1-\beta} = F_{ii}(1-\beta^T)$.
\end{proof}

\begin{corollary}[EWC as limiting case]
\label{cor:ewc_limit}
As $\beta \to 1$ with $T$ fixed, the EMA weight $(1-\beta)\beta^{T-t}$
converges to the uniform weight $1/T$ for each $t$.
Hence $v_i^{(T,\beta)} \to \frac{1}{T}\sum_{t=1}^T g_i(t)^2 = \hat{F}_{ii}^{\mathrm{unif}}$,
which is EWC's snapshot estimator applied over a full-training window.
\end{corollary}

\begin{remark}[Non-stationarity advantage]
\label{rem:nonstationarity}
Under a gradient process that decays from initial magnitude $M$ to stationary
magnitude $m \ll M$, EWC's uniform estimate on a window containing both phases
has non-stationarity bias $O((M-m)/N)$.
The EMA's exponential recency weighting yields bias $O(M\beta^T)$, which vanishes
exponentially and is strictly smaller for large $T$.
\end{remark}

\subsection{Elastic Penalty}
\label{sec:penalty}

After completing task $k$, \TCL{} freezes $v_i^{(T_k)}$ as the importance weight
for task $k$ and appends an elastic term to subsequent task losses:
\begin{equation}
  \mathcal{L}_{\TCL} = \mathcal{L}_{\mathrm{CE}}
    + \lambda \sum_{k<\text{current}} \sum_{i=1}^{d}
      v_i^{(T_k)} \cdot \bigl(\theta_i - \theta_i^{(k*)}\bigr)^2,
  \label{eq:tcl_penalty}
\end{equation}
where $\theta^{(k*)}$ is the parameter snapshot taken at the end of task $k$.
The penalty coefficient $\lambda$ is a global hyperparameter shared across all tasks.

\subsection{Governor Component (Non-Contributing)}
\label{sec:governor}

\TCL{} was originally designed with an additional \emph{governor} component:
a regime-detection mechanism that measures the ratio $\rho$ of current gradient
magnitude to a reference magnitude $\sigma^*$, and adjusts the learning rate
when $\rho$ falls below a threshold indicating convergence.
Our mechanism ablation (Section~\ref{sec:experiments}) demonstrates that:
(a) the governor-alone condition performs \emph{worse} than plain \SGD{}, and
(b) $\rho = 0$, $\sigma = 0$ in every production trace examined, indicating
the learning-rate adjustment has never fired.
Equation~\eqref{eq:tcl_penalty} represents the operative \TCL{} algorithm;
the governor is a dead branch.

\subsection{Connection to Fisher-EMA Theory}
\label{sec:theory_connection}

Equations \eqref{eq:ewc_snapshot} and \eqref{eq:ema_def} frame
the \TCL{}/\EWC{} comparison as a statistical estimation problem:
both methods estimate $F_{ii}$ from the same sequence of squared gradients,
but with different weighting schemes.
\EWC{} uses a uniform average over a terminal window;
\TCL{} uses an exponentially weighted average over the full training history.
Corollary~\ref{cor:ewc_limit} establishes \EWC{} as a special case,
and Theorem~\ref{thm:ema_convergence} and Remark~\ref{rem:nonstationarity}
provide the theoretical motivation for preferring the EMA under non-stationary
gradients.

% ── 4. Experiments ────────────────────────────────────────────────────────────
\section{Experiments}
\label{sec:experiments}

\subsection{Experimental Setup}
\label{sec:setup}

\paragraph{Dataset and protocol.}
We evaluate on Split-CIFAR-10~\cite{Krizhevsky2009} in the
task-incremental scenario: CIFAR-10 is partitioned into 5 tasks of 2 classes
each, with task identity provided at test time.
We use a ResNet-18~\cite{He2016} backbone with a separate classification head per
task.
All experiments use 40 training epochs per task, batch size 32,
and seeds $\{42, 0, 1, 2, 3\}$ ($n{=}5$ per condition).
This protocol is identical for all methods, removing architecture and
training-duration confounds.

\paragraph{Metrics.}
We report two primary metrics:
\emph{Catastrophic forgetting} $(\forgetmean{})$, defined as the average decrease
in per-task accuracy after all tasks are trained:
$\text{Forgetting} = \frac{1}{T{-}1}\sum_{k=1}^{T-1}(A_{k,k} - A_{T,k})$,
where $A_{i,j}$ is accuracy on task $j$ after training task $i$.
We also report final \emph{accuracy} $(\accmean{})$, the average accuracy across
all tasks after training completes, to detect degenerate low-forgetting solutions
that achieve low forgetting through accuracy collapse.

\paragraph{Statistical methodology.}
All confidence intervals use the $t$-distribution with $n{-}1{=}4$ degrees of
freedom ($t_{4,0.975} = 2.776$), replacing the common but incorrect $z = 1.96$
approximation that assumes $n \to \infty$.
Pairwise comparisons use paired one-tailed $t$-tests with Bonferroni family-wise
correction: for $k{=}4$ simultaneous comparisons within Phase 10
(TCL vs.\ EWC, SI, SGD, and an SI accuracy-collapse diagnostic),
the corrected significance threshold is $\alpha/k = 0.05/4 = 0.0125$.
Power is computed via the non-central $t$-distribution
($\mathrm{ncp} = -d\sqrt{n}$ for one-tailed tests).

\paragraph{Baselines.}
We compare against:
\textbf{\EWC{}} with $\lambda{=}100$ (the published default~\cite{Kirkpatrick2017}),
\textbf{\SI{}} with default hyperparameters~\cite{Zenke2017},
and plain \textbf{\SGD{}} (no regularization) as a lower bound.
DER++ and LwF are included in the pre-registered Phase 2 replication plan but
their GPU experiments have not yet run; we omit them from current results.

\subsection{Phase 10: Main Comparison (Split-CIFAR-10)}
\label{sec:phase10}

Table~\ref{tab:phase10} reports forgetting and accuracy for all four conditions
across $n{=}5$ seeds.
All confidence intervals use the $t$-distribution as described above.

\begin{table}[t]
  \caption{%
    Split-CIFAR-10 main comparison (Phase 10).
    $n{=}5$ seeds, $95\%$ CI via $t$-distribution ($t_{4,0.975}{=}2.776$).
    Bonferroni-corrected significance threshold $p < 0.0125$ ($k{=}4$).
    $\star$: Bonferroni-significant. \dag: directional ($p < 0.05$, uncorrected).
    Forgetting: lower is better. Accuracy: higher is better.
  }
  \label{tab:phase10}
  \centering
  \small
  \begin{tabular}{lcccc}
    \toprule
    & \multicolumn{2}{c}{Forgetting $\downarrow$} & \multicolumn{2}{c}{Accuracy $\uparrow$} \\
    \cmidrule(lr){2-3}\cmidrule(lr){4-5}
    Method & Mean & 95\% CI & Mean & 95\% CI \\
    \midrule
    \TCL{}             & 0.116 & \CI{0.080}{0.153} & 0.760 & \CI{0.729}{0.791} \\
    \EWC{} ($\lambda{=}100$) & 0.191 & \CI{0.162}{0.220} & 0.730 & \CI{0.702}{0.757} \\
    \SI{} (default)    & 0.092 & \CI{0.091}{0.093} & 0.500 & \CI{0.500}{0.500} \\
    \SGD{}             & 0.233 & \CI{0.229}{0.237} & 0.697 & \CI{0.693}{0.701} \\
    \bottomrule
  \end{tabular}
\end{table}

\begin{table}[t]
  \caption{%
    Pairwise comparisons versus \TCL{} (Phase 10, $n{=}5$, one-tailed paired $t$-test).
    $\Delta < 0$ means \TCL{} has lower (better) forgetting.
  }
  \label{tab:phase10_pairwise}
  \centering
  \small
  \begin{tabular}{lcccc}
    \toprule
    Comparison & $\Delta$ & $p$ (uncorr.) & Cohen's $d$ & Seeds \TCL{} better \\
    \midrule
    \TCL{} vs.\ \SGD{}  & $-0.117$ & $0.0012\,\star$ & $3.68$ & $5/5$ \\
    \TCL{} vs.\ \EWC{}  & $-0.075$ & $0.019\,\dag$  & $1.70$ & $5/5$ \\
    \TCL{} vs.\ \SI{}   & $+0.024$ & $0.142$        & $0.82$ & $1/5$ \\
    \bottomrule
    \multicolumn{5}{l}{%
      \footnotesize
      $\star$ Survives Bonferroni ($p < 0.0125$, $k{=}4$).
      \dag{} Directional ($p < 0.05$ uncorrected); does \emph{not} survive Bonferroni.
    }
  \end{tabular}
\end{table}

\paragraph{TCL vs.\ SGD.}
\TCL{} reduces forgetting by $0.117$ relative to \SGD{}
($p{=}0.0012$, Cohen's $d{=}3.68$, $5/5$ seeds, 100\% power at $n{=}5$).
This result survives Bonferroni correction and constitutes the single confirmed
headline finding of this paper.

\paragraph{TCL vs.\ EWC.}
\TCL{} reduces forgetting by $0.075$ relative to \EWC{} with default
hyperparameters ($p{=}0.019$, $d{=}1.70$, $5/5$ seeds, 92.5\% power at $n{=}5$).
This result does \textbf{not} survive Bonferroni correction ($p{=}0.019 > 0.0125$)
and must be treated as directional evidence.
At $d{=}1.70$ and $n{=}5$, the experiment had 92.5\% power, meaning the failure
to pass Bonferroni correction reflects the stringency of the family-wise threshold
rather than a weak effect.

\paragraph{SI accuracy collapse.}
\SI{} at published default hyperparameters achieves mean forgetting $0.092$,
\emph{lower} than \TCL{}'s $0.116$.
However, \SI{}'s accuracy of $0.500$ (chance level for a 2-class-per-task
benchmark) is constant across all 5 seeds with zero variance, indicating
catastrophic collapse rather than genuine retention.
A method that predicts randomly on every task has trivially low forgetting:
it was never accurate in the first place.
This illustrates why single-metric (forgetting-only) evaluation is insufficient,
and motivates reporting accuracy alongside forgetting.

\subsection{Phase 11: Mechanism Ablation}
\label{sec:phase11}

To isolate the productive component of \TCL{}, we ran a pre-registered
four-condition ablation: (1) full \TCL{} (EMA penalty + governor),
(2) penalty-only (EMA penalty, governor disabled),
(3) governor-only (LR adjustment, no penalty), and
(4) plain \SGD{} as a reference.
Results are in Table~\ref{tab:phase11}.

\begin{table}[t]
  \caption{%
    Mechanism ablation (Phase 11, Split-CIFAR-10).
    $n{=}5$ seeds.
    Pre-registered hypothesis: \emph{Does the governor contribute?}
    Verdict: governor-alone is worse than \SGD{}; penalty is the sole mechanism.
  }
  \label{tab:phase11}
  \centering
  \small
  \begin{tabular}{lcccc}
    \toprule
    & \multicolumn{2}{c}{Forgetting $\downarrow$} & \multicolumn{2}{c}{Accuracy $\uparrow$} \\
    \cmidrule(lr){2-3}\cmidrule(lr){4-5}
    Condition & Mean & 95\% CI & Mean & 95\% CI \\
    \midrule
    Full \TCL{}       & 0.127 & \CI{0.100}{0.154} & 0.766 & \CI{0.737}{0.795} \\
    Penalty-only      & 0.143 & \CI{0.121}{0.165} & 0.765 & \CI{0.748}{0.782} \\
    Governor-only     & 0.250 & \CI{0.165}{0.335} & 0.679 & \CI{0.606}{0.753} \\
    \SGD{}            & 0.219 & \CI{0.144}{0.295} & 0.709 & \CI{0.643}{0.775} \\
    \bottomrule
  \end{tabular}
\end{table}

\begin{table}[t]
  \caption{%
    Pairwise comparisons (Phase 11, full \TCL{} as reference, $n{=}5$).
    Pre-registered hypothesis was governor contribution; it is falsified.
  }
  \label{tab:phase11_pairwise}
  \centering
  \small
  \begin{tabular}{lcccc}
    \toprule
    Comparison & $\Delta$ & $p$ (uncorr.) & Cohen's $d$ & Seeds \TCL{} better \\
    \midrule
    Full vs.\ Governor-only & $-0.123$ & $0.023$ & $1.60$ & $5/5$ \\
    Full vs.\ \SGD{}        & $-0.092$ & $0.020$ & $1.68$ & $5/5$ \\
    Full vs.\ Penalty-only  & $-0.016$ & $0.315$ & $0.51$ & $4/5$ \\
    \bottomrule
    \multicolumn{5}{l}{%
      \footnotesize
      Full vs.\ penalty-only: $p{=}0.315$, underpowered ($24.5\%$ at $d{=}0.51$, $n{=}5$).
      Need $n{=}26$ for $80\%$ power. Null result is \emph{uninformative}.
    }
  \end{tabular}
\end{table}

\paragraph{Penalty is the mechanism.}
Governor-only performs \emph{worse} than \SGD{} ($0.250$ vs.\ $0.219$)
despite having the same objective, indicating that the regime-detection LR
adjustment actively interferes with learning when used without the penalty.
Full \TCL{} outperforms governor-only ($\Delta{=}{-}0.123$, $p{=}0.023$,
$d{=}1.60$, $5/5$ seeds), a result pre-registered as a falsification test of
the hypothesis ``governor contributes'' and classified as PUBLICATION\_ALLOWED
in our evidence inventory.
This negative finding---the governor mechanism is harmful without the penalty
and contributes nothing when combined with it---is the mechanistic conclusion of
this paper.

\paragraph{Governor never activates.}
Examination of all production traces confirms $\rho{=}0$, $\sigma{=}0$:
the gradient-ratio regime detector has never triggered the LR adjustment
in any evaluated configuration.
This is consistent with the governor-only ablation result.

\paragraph{Full vs.\ penalty-only.}
The comparison between full \TCL{} and penalty-only is not statistically resolved
($\Delta{=}{-}0.016$, $p{=}0.315$, $d{=}0.51$, $4/5$ seeds).
Power at $n{=}5$ for $d{=}0.51$ is only $24.5\%$, meaning the null result
is \emph{uninformative}: we cannot conclude the governor adds zero benefit at
$n{=}5$.
Resolving this comparison requires $n{=}26$ for $80\%$ power.
The available evidence supports that the penalty explains \emph{most} of the
\TCL{} benefit; the exact magnitude of any residual governor contribution
is undetermined.

% ── 5. Analysis ───────────────────────────────────────────────────────────────
\section{Analysis}
\label{sec:analysis}

\paragraph{EWC default hyperparameters underfit.}
\EWC{} with $\lambda{=}100$ achieves forgetting $0.191$, worse than \SGD{}'s
$0.233$ in the direction expected but with a gap that---directionally---favours
\TCL{}.
The Fisher-EMA theorem suggests the snapshot bias is structural:
for a ResNet-18 with 40 epochs per task, early-training gradients (epochs 1--10)
are dominated by loss curvature near random initialization, not by the task's
genuine parameter importance.
Mixing these into a uniform terminal window inflates importance weights for
parameters with high early-training sensitivity, over-constraining them in
subsequent tasks.
The EMA's exponential recency bias corrects this by down-weighting early batches
by $\beta^{T-t}$, focusing importance estimation on the convergence phase.

\paragraph{Effect size comparison.}
Cohen's $d{=}3.68$ for \TCL{} vs.\ \SGD{} is a very large effect by standard
conventions ($d > 0.8$: large; $d > 1.2$: very large).
Cohen's $d{=}1.70$ for \TCL{} vs.\ \EWC{} is also a large effect.
At $n{=}5$, power for $d{=}3.68$ is $100\%$ (confirmed headline result);
power for $d{=}1.70$ is $92.5\%$ (adequate, but $p{=}0.019$ fails
the Bonferroni-corrected threshold of $0.0125$).

\paragraph{Family-wise error control.}
Phase 10 tests four hypotheses simultaneously
(TCL vs.\ EWC, SI, SGD, and an SI accuracy diagnostic).
Without correction, the expected number of false positives at $\alpha{=}0.05$
is $0.2$ per family of four null hypotheses.
Bonferroni correction raises the per-comparison threshold to $0.0125$.
Only the TCL vs.\ SGD result ($p{=}0.0012$) passes this threshold.
The directional TCL vs.\ EWC finding motivates the pre-registered EWC replication
comparison (Phase 2, not yet run) as a pre-specified single-hypothesis test,
which would not require Bonferroni correction.

% ── 6. Limitations ────────────────────────────────────────────────────────────
\section{Limitations}
\label{sec:limitations}

\paragraph{1. The thermodynamic governor does not contribute.}
The regime-detection LR adjustment component of \TCL{}---the mechanism that
originally motivated the ``thermodynamic'' framing---is empirically falsified as
a productive mechanism.
Governor-alone produces \emph{higher} forgetting than plain \SGD{} ($0.250$
vs.\ $0.219$), and full \TCL{} outperforms governor-alone ($p{=}0.023$,
$d{=}1.60$, $5/5$ seeds, pre-registered falsification).
The benefit of the full \TCL{} system is attributable to the elastic penalty,
not to thermodynamic regime detection.

\paragraph{2. Regime-detection LR adjustment never fires.}
In every production trace examined, the gradient-magnitude ratio $\rho{=}0$
and the reference magnitude $\sigma{=}0$.
The learning-rate adjustment mechanism has never triggered in any evaluated
configuration.
This means all reported \TCL{} results were produced without the LR
adjustment, and the reported performance reflects the elastic penalty alone.

\paragraph{3. Task-incremental protocol only.}
All experiments provide task identity at test time (task-incremental scenario).
Results do not extend to class-incremental evaluation (where task identity is
withheld) without separate experimentation.
Class-incremental evaluation is substantially harder, and methods that reduce
forgetting in task-incremental settings often fail to generalize.

\paragraph{4. Single backbone (ResNet-18).}
All results are obtained with ResNet-18.
Generalization to Vision Transformers, wider residual networks, or other
architectures is untested.
The EMA importance accumulator is architecture-agnostic in principle, but
empirical validation requires separate evaluation.

\paragraph{5. TCL vs.\ EWC does not survive Bonferroni correction.}
The TCL vs.\ EWC comparison ($p{=}0.019$, $d{=}1.70$, $5/5$ seeds) does not
survive family-wise correction at $k{=}4$ (Bonferroni threshold $p{=}0.0125$).
This result is directional evidence only.
It motivates further pre-registered replication as a single-hypothesis test
(Phase 2, pre-registered but not yet run), which would not require Bonferroni
correction and would provide definitive confirmation or disconfirmation.

\paragraph{6. DER++ and LwF baselines not yet included.}
Dark Experience Replay (DER++) and Learning without Forgetting (LwF) are not
included in the current comparison.
Phase 2 GPU experiments (pre-registered, scripts written) include both baselines
with fair hyperparameter selection, but those experiments have not yet been
executed.
Current results should therefore not be interpreted as establishing \TCL{}'s
position in a complete benchmark; they establish the EMA mechanism finding and
the SGD gap specifically.

% ── 7. Conclusion ─────────────────────────────────────────────────────────────
\section{Conclusion}

We have presented a formal and empirical argument that \EWC{}'s terminal-snapshot
importance estimator is suboptimal under non-stationary training gradients,
and that replacing it with a continuous exponential moving average---the
ThermalImportance accumulator---corrects this bias.
The resulting method, \TCL{}, reduces catastrophic forgetting on Split-CIFAR-10
by $0.117$ relative to \SGD{} ($p{=}0.0012$, $d{=}3.68$, Bonferroni-corrected),
and shows a large directional improvement over \EWC{} with default hyperparameters
($\Delta{=}{-}0.075$, $d{=}1.70$, $5/5$ seeds, directional pending replication).

A pre-registered mechanism ablation definitively falsifies the hypothesis that
\TCL{}'s thermodynamic governor component contributes to performance.
Governor-alone performs worse than \SGD{}; the elastic penalty is the entire
operative mechanism.

The honest scope of these results is narrow: task-incremental only, ResNet-18 only,
Split-CIFAR-10 only.
The confirmed headline finding is the SGD gap.
The EWC gap requires pre-registered replication (Phase 2 scripts written and ready).
We report this work at workshop stage with the intention of completing the
replication before any archival submission.

\section*{Acknowledgements}

Omitted for anonymous review.

% ── bibliography ──────────────────────────────────────────────────────────────
\bibliographystyle{plainnat}
\bibliography{paper}

\end{document}
